\chapter{Szybka sciaga skladni}

Ten dodatek nie zastepuje pelnych rozdzialow. Ma sluzyc jako szybkie przypomnienie
najwazniejszych konstrukcji.

\section{Podstawowy szkielet}
\begin{lstlisting}[style=moonchunk]
chunk "Main" {
  output: "./dist";

  page "" {
    content {
      <h1>Hello</h1>
    };
  };
};

moon(Main);
\end{lstlisting}

\section{Zmienne i stale}
\begin{lstlisting}[style=moonchunk]
let a: int = 1;
let b: string;
b = "hello";
const title: string = "MoonChunk";
\end{lstlisting}

\section{Warunki}
\begin{lstlisting}[style=moonchunk]
if (a > 0) {
  print("plus");
} else {
  print("zero lub minus");
};
\end{lstlisting}

\section{Petle}
\begin{lstlisting}[style=moonchunk]
for (let i: int = 0; i < 10; i++) {
  print(i);
};

while (a < 5) {
  a = a + 1;
};
\end{lstlisting}

\section{Funkcje}
\begin{lstlisting}[style=moonchunk]
function add(a: int, b: int): int {
  return a + b;
}

sum(a: int, b: int): int => a + b;
\end{lstlisting}

\section{Importy}
\begin{lstlisting}[style=moonchunk]
import { Metadata } from "./metadata.mncnk";
import * as Fragments from "./fragments.mncnk";
\end{lstlisting}

\section{Include}
\begin{lstlisting}[style=moonchunk]
chunk "Main" {
  @include Metadata;
  @include Fragments.Hero;
  output: "./dist";
};
\end{lstlisting}

\section{Tablice i obiekty}
\begin{lstlisting}[style=moonchunk]
const arr = [1, 2, 3];
const obj = { title: "Hello", views: 10 };
print(arr[0]);
print(obj.title);
\end{lstlisting}

\section{Rzutowanie}
\begin{lstlisting}[style=moonchunk]
let a: int = "42" as int;
let b: string = 10 as string;
let c: bool = (bool)3;
let d: bool = 3 as unknown as bool;
\end{lstlisting}

\section{Globale}
\begin{lstlisting}[style=moonchunk]
env {
  global siteName: string = "MoonChunk";
};
\end{lstlisting}
